\documentclass[10pt,a4paper]{amsart}
\usepackage{amsmath,amssymb,amsthm,graphicx}
% Retained minimal preamble for the Stacks Project chapter more-algebra.tex (More on Algebra).
% The chapter file pinned-original-source.tex carries only \input{preamble} on its first line, so no
% preamble text lives inside the chapter itself.  The Stacks Project uses one single shared preamble for
% every chapter; the snapshot stages it identically as sources/stacks-chow/preamble.tex and
% sources/stacks-etale-cohomology/preamble.tex.  The block below is copied mechanically, line for line,
% from sources/stacks-chow/preamble.tex lines 14-18 (the Xy-pic setup) and 154-217 (the native theorem
% declarations together with the mathematical macros).  Nothing was added, removed or redefined; the
% wrapper class and the amsmath/amssymb/amsthm/graphicx base packages are supplied by the assembler.
% Chapter: more-algebra.tex -- More on Algebra.
% For commutative diagrams we use Xy-pic
\usepackage[all]{xy}

% We use 2cell for 2-commutative diagrams.
\xyoption{2cell}

% Theorem environments.
%
\theoremstyle{plain}
\newtheorem{theorem}[subsection]{Theorem}
\newtheorem{proposition}[subsection]{Proposition}
\newtheorem{lemma}[subsection]{Lemma}

\theoremstyle{definition}
\newtheorem{definition}[subsection]{Definition}
\newtheorem{example}[subsection]{Example}
\newtheorem{exercise}[subsection]{Exercise}
\newtheorem{situation}[subsection]{Situation}

\theoremstyle{remark}
\newtheorem{remark}[subsection]{Remark}
\newtheorem{remarks}[subsection]{Remarks}

\numberwithin{equation}{subsection}

% Macros
%
\def\lim{\mathop{\mathrm{lim}}\nolimits}
\def\colim{\mathop{\mathrm{colim}}\nolimits}
\def\Spec{\mathop{\mathrm{Spec}}}
\def\Hom{\mathop{\mathrm{Hom}}\nolimits}
\def\Ext{\mathop{\mathrm{Ext}}\nolimits}
\def\SheafHom{\mathop{\mathcal{H}\!\mathit{om}}\nolimits}
\def\SheafExt{\mathop{\mathcal{E}\!\mathit{xt}}\nolimits}
\def\Sch{\mathit{Sch}}
\def\Mor{\mathop{\mathrm{Mor}}\nolimits}
\def\Ob{\mathop{\mathrm{Ob}}\nolimits}
\def\Sh{\mathop{\mathit{Sh}}\nolimits}
\def\NL{\mathop{N\!L}\nolimits}
\def\CH{\mathop{\mathrm{CH}}\nolimits}
\def\proetale{{pro\text{-}\acute{e}tale}}
\def\etale{{\acute{e}tale}}
\def\QCoh{\mathit{QCoh}}
\def\Ker{\mathop{\mathrm{Ker}}}
\def\Im{\mathop{\mathrm{Im}}}
\def\Coker{\mathop{\mathrm{Coker}}}
\def\Coim{\mathop{\mathrm{Coim}}}

% Boxtimes
%
\DeclareMathSymbol{\boxtimes}{\mathbin}{AMSa}{"02}

%
% Macros for moduli stacks/spaces
%
\def\QCohstack{\mathcal{QC}\!\mathit{oh}}
\def\Cohstack{\mathcal{C}\!\mathit{oh}}
\def\Spacesstack{\mathcal{S}\!\mathit{paces}}
\def\Quotfunctor{\mathrm{Quot}}
\def\Hilbfunctor{\mathrm{Hilb}}
\def\Curvesstack{\mathcal{C}\!\mathit{urves}}
\def\Polarizedstack{\mathcal{P}\!\mathit{olarized}}
\def\Complexesstack{\mathcal{C}\!\mathit{omplexes}}
% \Pic is the operator that assigns to X its picard group, usage \Pic(X)
% \Picardstack_{X/B} denotes the Picard stack of X over B
% \Picardfunctor_{X/B} denotes the Picard functor of X over B
\def\Pic{\mathop{\mathrm{Pic}}\nolimits}
\def\Picardstack{\mathcal{P}\!\mathit{ic}}
\def\Picardfunctor{\mathrm{Pic}}
\def\Deformationcategory{\mathcal{D}\!\mathit{ef}}

% --- appended 2026-10-12 (Chief): the project own definitions of the reference/slogan/history
% display environments, copied mechanically from the shared Stacks preamble lines 8-12
% (\usepackage{verbatim} plus the three \newenvironment lines, which the chapter preamble keeps
% verbatim from the project preamble). Nothing added, removed or redefined: these environments
% carry no mathematical content and the project itself typesets them as comments. Without these
% three lines a counted statement whose own text uses them fails with Environment slogan undefined.
\usepackage{verbatim}
\newenvironment{reference}{\comment}{\endcomment}
\newenvironment{slogan}{\comment}{\endcomment}
\newenvironment{history}{\comment}{\endcomment}
\usepackage{cleveref}
\crefname{originallemma}{original lemma}{original lemmas}
\Crefname{originallemma}{Original lemma}{Original lemmas}
\makeatletter
\newlabel{lemma-projective-amplitude}{{0A5P}{source}{}{sourcepointer-1}{}}
\newlabel{lemma-projective-amplitude@cref}{{[originallemma][0][]0A5P}{[1][0][]source}{}{}{}}
\makeatother
\crefname{originallemma}{original lemma}{original lemmas}
\Crefname{originallemma}{Original lemma}{Original lemmas}
\makeatletter
\newlabel{lemma-Noetherian-pseudo-coherent}{{066E}{source}{}{sourcepointer-2}{}}
\newlabel{lemma-Noetherian-pseudo-coherent@cref}{{[originallemma][0][]066E}{[1][0][]source}{}{}{}}
\makeatother
\crefname{originallemma}{original lemma}{original lemmas}
\Crefname{originallemma}{Original lemma}{Original lemmas}
\makeatletter
\newlabel{lemma-projective-amplitude-pseudo-coherent}{{0G9A}{source}{}{sourcepointer-3}{}}
\newlabel{lemma-projective-amplitude-pseudo-coherent@cref}{{[originallemma][0][]0G9A}{[1][0][]source}{}{}{}}
\makeatother
\title{Vanishing of Ext\textasciicircum{}1 against all modules versus two-term projective complexes, with the Noetherian finite variant}
\author{The Stacks Project Authors}
\date{stacks-project:more-algebra.tex:local-snapshot-20261012; GNU FDL 1.2 or later}
\begin{document}
\maketitle
\setcounter{subsection}{1}
\noindent\textit{Editorial note.} All delivered mathematical text is the author's own, copied from the pinned Stacks Project chapter file more-algebra.tex; no statement or proof step was rewritten, reordered or removed. The counted span is the author environment and its own proof environment. References inside the span are handed to the assembler as navigation stubs that carry the canonical tag of the cited tag only; whether each of them is load-bearing is recorded on the stub itself.\par\medskip
\noindent\textit{Allowed standard prerequisites (editorial).}\begin{itemize}
\item A complex with cohomology concentrated in [-1, 0] has vanishing Ext\textasciicircum{}1 against every module if and only if H\textasciicircum{}\{-1\} = 0 and H\textasciicircum{}0 is projective. invoked at proof lines 22280-22282. Gives the equivalence of (1) and (2).
\item A complex over a Noetherian ring with finite cohomology modules in degrees -1 and 0 is pseudo-coherent. invoked at proof lines 22283-22285. Lets the pseudo-coherent version of the criterion apply.
\item For a pseudo-coherent complex the Ext\textasciicircum{}1 criterion may be tested against finite modules only. invoked at proof lines 22285-22287. Gives the equivalence with (3).
\end{itemize}\medskip
\ifdefined\hypertarget\hypertarget{sourcepointer-1}{}\fi
\noindent\textit{Source reference.} Original lemma 0A5P: sources/stacks-more-algebra/pinned-original-source.tex, source lines 17414--17414. Cited at proof lines 22280-22281 for the equivalence of (1) and (2); the conclusion is used inside the counted span.\par\medskip
\ifdefined\hypertarget\hypertarget{sourcepointer-2}{}\fi
\noindent\textit{Source reference.} Original lemma 066E: sources/stacks-more-algebra/pinned-original-source.tex, source lines 16517--16517. Cited at proof lines 22283-22284 for the pseudo-coherence of K; the conclusion is used inside the counted span.\par\medskip
\ifdefined\hypertarget\hypertarget{sourcepointer-3}{}\fi
\noindent\textit{Source reference.} Original lemma 0G9A: sources/stacks-more-algebra/pinned-original-source.tex, source lines 20136--20136. Cited at proof lines 22285-22286 for the equivalence with (3); the conclusion is used inside the counted span.\par\medskip
\begin{lemma}
\label{lemma-ext-1-zero}
Let $R$ be a ring. Let $K \in D(R)$ with $H^i(K) = 0$ for
$i \not \in \{-1, 0\}$. The following are equivalent
\begin{enumerate}
\item $H^{-1}(K) = 0$ and $H^0(K)$ is a projective module and
\item $\Ext^1_R(K, M) = 0$ for every $R$-module $M$.
\end{enumerate}
If $R$ is Noetherian and $H^i(K)$ is a finite $R$-module for
$i = -1, 0$, then these are also equivalent to
\begin{enumerate}
\item[(3)] $\Ext^1_R(K, M) = 0$ for every finite $R$-module $M$.
\end{enumerate}
\end{lemma}


\begin{proof}
The equivalence of (1) and (2) follows from
Lemma \ref{lemma-projective-amplitude}.
If $R$ is Noetherian and $H^i(K)$ is a finite $R$-module for
$i = -1, 0$, then $K$ is pseudo-coherent, see
Lemma \ref{lemma-Noetherian-pseudo-coherent}.
Thus the equivalence of (1) and (3) follows from
Lemma \ref{lemma-projective-amplitude-pseudo-coherent}.
\end{proof}

\end{document}
